\chapter{Исследовательская часть}
% в исследовательском разделе отобразить полученные результаты для приложенного набора файлов, оценить сложность полученного силами БЯМ алгоритма решения задачи.

\section{Тестирование}

\begin{table}[!htbp]
	\caption{Результаты тестирования.}
	\centering
	\begin{tabularx}{\textwidth}{|*{3}{>{\centering\arraybackslash}X|}}
		\hline
		\textbf{Модель} & \textbf{Вход} & \textbf{Выход} \\\hline
		qwen3-coder & Отчёт по первой лабораторной работе & Случаи использования 'оглавление' вместо 'содержание' не найдены \\\hline
		devstral-2  & Отчёт по первой лабораторной работе & В файле используется 'содержание' (не 'оглавление'). \\\hline
		cogito-2.1  & Отчёт по первой лабораторной работе & Упоминания 'оглавление' или 'содержание' не найдены. \\\hline
		gemma3      & Отчёт по первой лабораторной работе & Обнаружено содержание.\newlineОглавление не найдено. \\\hline
	\end{tabularx}
\end{table}

\section{Трудоёмкость}

С помощью введённой модели вычислений (табл.~\ref{t1}) была определена трудоёмкость ранее полученных реализаций, результаты представлены в табл.~\ref{t2}.

\begin{table}[!htbp]
	\caption{Цены базовых операций.}
	\label{t1}
	\centering
	\begin{tabularx}{\textwidth}{|*{2}{>{\centering\arraybackslash}X|}}
		\hline
		\textbf{Операции} & \textbf{Цена} \\\hline
		{ +, - , =, ++, $--$ , <<, >> } & 1 \\\hline
		{ *, /, \%, *=, /= } & 2 \\\hline
		{ ==, >=, <=, !=, >, <, \&\&, || } & 1 \\\hline
		{ [], \{\} } & 1\\\hline
	\end{tabularx}
\end{table}
\FloatBarrier

\begin{table}[!htbp]
	\centering
	\caption{Трудоёмкость алгоритмов.}
	\label{t2}
	\small
	\begin{tabularx}{\textwidth}{|p{3cm}|*{2}{X|}}
		\hline
		\textbf{БЯМ} & \textbf{Трудоёмкость ЛС} & \textbf{Трудоёмкость ХС} \\\hline
		qwen3-coder & $90n + 342$ & $O(60n)$ \\\hline
		devstral-2 & $11n + 45$ & $O(n)$ \\\hline
		cogito-2.1 & $37n + 68$ & $O(23n)$ \\\hline
		gemma3 & $8n$ & $O(10n)$ \\\hline
	\end{tabularx}
\end{table}
\FloatBarrier

\section*{Вывод}
%<что сделали и получили в результате, кратко>

В данном разделе было проведено тестирование программ, написанных БЯМ, и рассчитана их трудоёмкость.

\clearpage
